\subsection{Weighted Einstein--Vlasov Particle--Mesh Module}
\label{sec:weighted-einstein-vlasov-pic}

This collisionless module evolves a spherical relativistic marker ensemble
and its self-consistent spacetime in Schwarzschild coordinates.  Separate task
executables select a low-compactness halo or high-compactness collapse, so a
parameter file cannot silently switch between those physical problems.  The
module contains no SIDM collision operator and does not evolve an expanding
cosmological background.  Schwarzschild coordinates approach a coordinate
singularity as \(2m/r\rightarrow1\); the module therefore diagnoses horizon
approach but is not a horizon-penetrating black-hole evolution scheme.

\noindent\textbf{Maturity and validation.}
The task-specific runners implement the Schwarzschild weighted
particle-in-cell (WPIC) method.  The low-compactness task compares
Einstein--Vlasov and Vlasov--Poisson evolution from the same sampled initial
positions and momenta.  Each module assigns the weights required by its own
phase-space measure.  Quantitative collapse results require joint refinement
of marker number, radial mesh, timestep, omitted-weight control, and sampling
realization, together with constraint, mass, binding, and outcome checks.
Horizon-penetrating evolution belongs to the EF module.

\begin{figure}[htbp]
\centering
\textbf{Physics flow: Einstein--Vlasov WPIC}
\par\smallskip
\begin{tikzpicture}[node distance=0.22cm]
\node[physicsflow] (ew1) {Create weighted spherical markers carrying position, radial momentum, and angular momentum};
\node[physicsflow,below=of ew1] (ew2) {Deposit their rest mass, relativistic energy, pressures, and radial current onto radial shells};
\node[physicsflow,below=of ew2] (ew3) {Solve the spherical Einstein constraints to obtain enclosed mass, radial metric, and lapse};
\node[physicsflow,below=of ew3] (ew4) {Predict every marker to the middle of the timestep using the old spacetime};
\node[physicsflow,below=of ew4] (ew5) {Deposit the predicted markers and solve the midpoint spacetime};
\node[physicsflow,below=of ew5] (ew6) {Move the original markers through the full timestep using that midpoint spacetime};
\node[physicsflow,below=of ew6] (ew7) {Recompute the endpoint matter and spacetime, reducing the timestep if the coordinate patch would fail};
\node[physicsflow,below=of ew7] (ew8) {Write profiles, conservation diagnostics, and an atomic restart; then continue};
\draw[physicsarrow] (ew1) -- (ew2);
\draw[physicsarrow] (ew2) -- (ew3);
\draw[physicsarrow] (ew3) -- (ew4);
\draw[physicsarrow] (ew4) -- (ew5);
\draw[physicsarrow] (ew5) -- (ew6);
\draw[physicsarrow] (ew6) -- (ew7);
\draw[physicsarrow] (ew7) -- (ew8);
\end{tikzpicture}
\end{figure}

\subsubsection{Physical variables and invariant marker weight}

The Schwarzschild line element is
\[
 \dd s^2=-e^{2\mu(\tau,r)}\dd\tau^2
          +e^{2\lambda(\tau,r)}\dd r^2+r^2\dd\Omega^2,    \tag{EWP1}
\]
where \(r\) is areal radius, \(\tau\) is Schwarzschild coordinate time,
\(\mu\) is the lapse
potential, \(\lambda\) is the radial metric potential, and
\(\dd\Omega^2\) is the unit-sphere metric.  In the source,
\code{WeightedPICMarker} is the record for one persistent numerical sample of
the distribution.  A marker stores
\[
 (m_p,r_p,w_p,\eta_p,I_p),\qquad
 \eta_p=\frac{L_p^2}{2},\qquad
 \epsilon_p=\sqrt{1+w_p^2+\frac{2\eta_p}{r_p^2}}.        \tag{EWP2}
\]
Here \(m_p\) is invariant rest-number weight, \(w_p\) is orthonormal radial
momentum measured by a local static observer, \(L_p\) is the conserved
angular-momentum magnitude, \(I_p\) is a persistent identifier, and
\(\epsilon_p\) is energy per unit rest mass in that local orthonormal frame.

Because \(w\) is orthonormal rather than canonical, the invariant reduced
phase measure is
\[
 \dd\Gamma_{\rm inv}
 =8\pi^2 e^\lambda\,\dd r\,\dd w\,\dd\eta.               \tag{EWP3}
\]
An initial sampler first constructs coordinate weights using
\(\dd r\,\dd w\,\dd\eta\), then iterates
\[
 m_p^{(n+1)}
 =m_{p,\rm coord}\exp[\lambda^{(n)}(r_p)]                \tag{EWP4}
\]
until the relative change is below \(2\times10^{-12}\).  This fixed-point
iteration makes \(m_p\) the invariant weight required by equation (EWP3).
This metric factor converts coordinate-volume sampling into invariant
particle number.  Omitting it changes the matter source and produces a false
growth of Arnowitt--Deser--Misner (ADM) mass, the total gravitating mass read
from the asymptotic exterior metric.

There is one explicit interface exception.  Markers returned by
\path{cpu/include/ev/weighted_einstein_vlasov_static_shell.hpp} already carry the
fully invariant quadrature weight.  For its continuum distribution
\(f=C(1-E/E_0)_+^k\delta(L-L_0)\), \(C\) is the distribution amplitude,
\(L_0>0\) is the single angular-momentum magnitude carried by the shell,
\((x)_+=\max(x,0)\) keeps only the positive part, and \(\delta\) is the Dirac
delta distribution that enforces \(L=L_0\).
The cutoff power satisfies \(1\le k\le8\), with default \(k=1\).
\(E_p=e^{\mu(r_p)}\epsilon_p\) is marker \(p\)'s conserved energy in the
static metric, and \(E_0\) is the upper energy cutoff.  The factors
\(\Delta r_p\) and \(\Delta w_p\) are the radial and positive-radial-momentum
quadrature measures assigned to that sample; each includes its numerical
quadrature weight rather than denoting a universal mesh spacing.  The returned
marker weight is
\[
 m_p=8\pi^2 CL_0 e^{\lambda(r_p)}
      \left(1-\frac{E_p}{E_0}\right)^k
      \Delta r_p\,\Delta w_p .
\]
The \code{energy_cutoff_shell} driver therefore passes those weights directly
to the Schwarzschild WPIC solver and must \emph{not} call
\code{canonicalize_weighted_einstein_vlasov_pic_initial_weights};
doing so would insert \(e^\lambda\) twice.  Its exact continuum equations,
delta-function normalization, sampling rules, fixed inward perturbation, and
threshold interface are documented in
Section~\ref{sec:ef-static-energy-cutoff-shell}.

\subsubsection{Relativistic deposition}

Markers deposit into shell centers placed uniformly in the volume coordinate
\[
 q=\frac{r^3}{3},\qquad \Delta V_i=4\pi\Delta q_i.       \tag{EWP5}
\]
CIC means cloud in cell, a two-center linear particle shape; TSC means
triangular shaped cloud, a smoother three-center quadratic shape.  TSC is the
default.  Let \(a_{pi}\) be the nonnegative shape weight of marker \(p\) in
shell \(i\), with \(\sum_i a_{pi}=1\).  The deposited shell-extensive
quantities are
\[
\begin{aligned}
 {\cal M}_{0,i}&=\sum_p a_{pi}m_p,\\
 {\cal E}_i&=\sum_p a_{pi}m_p\epsilon_p,\\
 {\cal P}_{r,i}&=\sum_p a_{pi}m_p\frac{w_p^2}{\epsilon_p},\\
 {\cal P}_{t,i}&=\sum_p a_{pi}m_p
                 \frac{2\eta_p/r_p^2}{\epsilon_p},\\
 {\cal J}_i&=\sum_p a_{pi}m_pw_p.
\end{aligned}                                           \tag{EWP6}
\]
\({\cal M}_0\) is rest mass, \({\cal E}\) is the invariant-weighted local
particle energy, the two
\({\cal P}\)'s generate radial and tangential pressure, and \({\cal J}\) is
radial current.  Unequal-weight sampling quality is reported by
\[
 N_{{\rm eff},i}
 =\frac{(\sum_p a_{pi}m_p)^2}{\sum_p(a_{pi}m_p)^2}.      \tag{EWP7}
\]

\subsubsection{Einstein constraint solve}

Define the metric gap
\[
 g(r)=1-\frac{2m(r)}{r}=e^{-2\lambda(r)},                \tag{EWP8}
\]
where \(m(r)\) is the Misner--Sharp enclosed mass.  Inside one shell, the
source \({\cal E}_i/\Delta q_i\) is held constant and the radial mass
constraint is integrated directly:
\[
 \frac{\dd m}{\dd q}
 =\frac{{\cal E}_i}{\Delta q_i}
  \sqrt{1-\frac{2m}{r(q)}}.                              \tag{EWP9}
\]
The square root is the \(e^{-\lambda}\) conversion from invariant canonical
marker measure to physical stress energy.  VENUS uses an implicit-midpoint
step with step doubling.  A radial segment is never accepted if its outer
face has \(2m/r\ge1\); this is a coordinate-patch condition, not an artificial
density cap.

For a radial segment \([q_L,q_R]\), write \(m_L=m(q_L)\),
\(r_m=[3(q_L+q_R)/2]^{1/3}\), and
\(Q=({\cal E}_i/\Delta q_i)(q_R-q_L)\).  The implicit-midpoint equation
for the mass increment \(\delta m=m_R-m_L\) is
\[
 \delta m=Q\sqrt{g_L-\delta m/r_m},\qquad
 g_L=1-2m_L/r_m.
\]
Squaring gives \((\delta m)^2+(Q^2/r_m)\delta m-Q^2g_L=0\).
The nonnegative root is evaluated in the cancellation-resistant form
\[
 \delta m=\frac{2Q^2g_L}
 {\sqrt{(Q^2/r_m)^2+4Q^2g_L}+Q^2/r_m}.       \tag{EWP9a}
\]
For zero source the mass stays \(m_L\).  The source implements this root
in \code{radial_mass_midpoint_step}:
\codesource{Schwarzschild WPIC implicit mass increment}
{\path{cpu/include/ev/weighted_einstein_vlasov_pic_schwarzschild.hpp}}
{Evaluates the positive quadratic root of the radial midpoint constraint}
\begin{lstlisting}[language=C++]
const long double physical_energy =
    denominator > 0.0L
    ? 2.0L * canonical_energy_squared * gap_before_step
        / denominator
    : 0.0L;
result.mass = mass_left + physical_energy;
\end{lstlisting}
Here \code{canonical_energy_squared} is \(Q^2\),
\code{gap_before_step} is \(g_L\), and \code{denominator} is the
denominator in (EWP9a).  One full segment and two half segments estimate the
radial integration error.  Refinement continues until that error and the
positive outer-face gap satisfy the solver's acceptance conditions.

At a shell center the physical moments are
\[
 \rho_i=\frac{{\cal E}_i\sqrt{g_i}}{\Delta V_i},\quad
 p_i=\frac{{\cal P}_{r,i}\sqrt{g_i}}{\Delta V_i},\quad
 p_{t,i}=\frac{{\cal P}_{t,i}\sqrt{g_i}}{\Delta V_i},\quad
 j_i=\frac{{\cal J}_i\sqrt{g_i}}{\Delta V_i}.            \tag{EWP10}
\]
The remaining radial constraints are
\[
\lambda=-\frac12\ln g,\qquad
\lambda_r=\frac{4\pi r\rho-m/r^2}{g},\qquad
\mu_r=\frac{m/r^2+4\pi rp}{g}.                           \tag{EWP11}
\]
The outer lapse equals the exterior Schwarzschild value,
\[
 \mu(R)=\frac12\ln g(R),                                 \tag{EWP12}
\]
and \(\mu_r\) is integrated inward.  This matching assumes vacuum beyond
the outer field face \(R\), with Schwarzschild time normalized at infinity.
Matter outside \(R\) contributes to that lapse normalization even when its
mass lies outside an interior particle's enclosed-mass force.  The production
low- and high-compactness calculations therefore require all positive-weight
matter to remain inside the vacuum boundary.  The Schwarzschild donor uses
the same condition when constructing null initial data.

The self-consistent update checks the initial, predicted midpoint, and trial
endpoint marker states before committing the particle array.  A trial that
escapes with positive rest weight leaves the accepted state unchanged.
Production enables this condition explicitly:
\par\noindent\begin{minipage}{\linewidth}
\codesource{Schwarzschild isolated outer-boundary guard}
{\path{cpu/src/drivers/einstein_vlasov_wpic.cpp};
 \path{cpu/include/ev/weighted_einstein_vlasov_pic_schwarzschild.hpp}}
{Enforces the vacuum exterior before production evolution and rejects a trial
 that would remove matter through the outer field face}
\begin{lstlisting}[language=C++]
step_options.require_vacuum_outer_boundary = true;
ev::check_weighted_einstein_vlasov_pic_vacuum_outer_boundary(markers);
\end{lstlisting}
\end{minipage}\par
The driver then saves the last accepted checkpoint and its history/profile
endpoint.  The high-compactness outcome is
\code{sc_outer_boundary_evolution_limit}; its summary records accepted time,
rejected trial duration, and the reason under
\code{schwarzschild_evolution_status}.  This limited interval remains an
incomplete physical evolution.  A larger vacuum domain supplies the
corresponding continuation problem.  The library default retains finite-domain
outflow bookkeeping for callers that provide the exterior matter and lapse
model; these escaped markers keep their invariant weights.

The time derivative needed by the
characteristics follows from the radial current:
\[
 \lambda_\tau=-4\pi r e^\mu\,j_{\rm canonical},\qquad
 j_{\rm canonical}=\frac{{\cal J}}{\Delta V}.            \tag{EWP13}
\]
The implementation reconstructs the lapse and its derivative from cell
centers and faces.  Monotonic slope limiting keeps the interpolation
consistent with the attractive constraint force at those nodes, so sampling
noise cannot create an artificial repulsive interval between them.

\subsubsection{Relativistic marker characteristics}

With \(L=\sqrt{2\eta}\), a collisionless marker follows
\[
 \frac{\dd r}{\dd\tau}
 =e^{\mu-\lambda}\frac{w}{\epsilon},                    \tag{EWP14}
\]
\[
 \frac{\dd w}{\dd\tau}
 =-\lambda_\tau w
  -e^{\mu-\lambda}\mu_r\epsilon
  +e^{\mu-\lambda}\frac{L^2}{r^3\epsilon}.              \tag{EWP15}
\]
The three terms in equation (EWP15) are the time-dependent radial-metric
term, gravitational/lapse force, and angular-momentum barrier.  Rest weight
\(m_p\) and \(\eta_p\) remain invariant.

An ordinary fourth-order Runge--Kutta step advances most markers.  Near a true
origin, the code rewrites the motion using two Cartesian components, so the
apparently singular \(L^2/r^3\) term becomes ordinary smooth central-force
motion.  For a resolved nonzero inner boundary, rapidly turning low-pericentre
orbits can use the Levi--Civita regularized characteristic integrator.  Here
``regularized'' means that the independent variable and coordinates are
changed so a close passage is numerically smooth; the physical orbit is not
softened.

\subsubsection{Self-consistent midpoint time step}

One accepted interval \(\Delta\tau\) performs:
\begin{enumerate}
\item deposit markers at \(\tau^n\) and solve the old constraints;
\item push a copy by \(\Delta\tau/2\) in the old metric;
\item deposit the copy and solve the midpoint constraints;
\item push the original markers by \(\Delta\tau\) in the midpoint metric;
\item deposit the endpoint markers and solve the endpoint constraints.
\end{enumerate}
This is a second-order self-consistent midpoint coupling even though each
frozen-metric characteristic uses RK4 internally.  If the predictor or
endpoint leaves the Schwarzschild patch, the complete trial is discarded and
\(\Delta\tau\) is bisected, up to the configured bisection cap.  The
\code{gap_stop} value is a requested Schwarzschild coordinate-patch stopping
proxy; it is not substituted into the field equations and is neither an
apparent-horizon finder nor a certified black-hole event.

The predictor starts from a copy, while the full push starts again from the
original ensemble.  The following selected source lines show those time
levels; the intervening patch-validity checks are omitted:
\codesource{Schwarzschild WPIC midpoint predictor and full push}
{\path{cpu/include/ev/weighted_einstein_vlasov_pic_schwarzschild.hpp}}
{Uses the old metric for the half-step prediction and the redeposited midpoint metric for the full step}
\begin{lstlisting}[language=C++,basicstyle=\ttfamily\scriptsize]
std::vector<WeightedPICMarker> midpoint_markers = markers;
advance_weighted_einstein_vlasov_pic_frozen_metric_parallel(
    midpoint_markers, 0.5 * trial_dt, grid, old_field.metric,
    options.threads, options.characteristic_substeps);
// After depositing and solving midpoint_markers:
endpoint_markers = markers;
advance_weighted_einstein_vlasov_pic_frozen_metric_parallel(
    endpoint_markers, trial_dt, grid, midpoint_field.metric,
    options.threads, options.characteristic_substeps);
\end{lstlisting}

The timestep scans marker mesh-crossing and momentum-change times.  The
omitted-mass control listed below allows the fastest weighted tail, up to that
fraction of total rest mass, to be omitted from the global step estimate.
This is a weighted quantile, not a raw marker fraction.
It prevents a negligibly weighted importance-sampled marker from controlling
every step, but it must be reduced in the convergence study.

\subsubsection{Initial-condition provenance and acceptance checks}

There are two consuming tasks:
\begin{itemize}
\item \code{venus_cpu_einstein_vlasov_wpic_low_compactness}, selected by
\path{cpu/src/apps/einstein_vlasov/wpic_low_compactness.cpp};
\item \code{venus_cpu_einstein_vlasov_wpic_black_hole}, selected by
\path{cpu/src/apps/einstein_vlasov/wpic_black_hole.cpp}.
\end{itemize}
The shared generator, importer, and evolution driver is
\path{cpu/src/drivers/einstein_vlasov_wpic.cpp}.

The low task generates one deterministic Plummer marker vector and copies the
same phase points into the relativistic and Newtonian states.  The high task
either generates Plummer or Gaussian-shell markers, imports an immutable
physical-initial-condition file, or invokes the bound energy-cutoff continuum
and sampler in
\path{cpu/include/ev/weighted_einstein_vlasov_static_shell.hpp}.

The generated dynamical object is an in-memory
\code{std::vector<WeightedPICMarker>}, a resizable list of the marker records
defined above.  With \code{--physical-ic PATH}, the high task also writes or
authenticates a version-1 \code{VENUSEVPHYSIC1} binary object at that path.
This is an immutable physical initial condition, not a time-evolution
restart.  Its payload is, in order: a 16-byte format tag, a 32-bit version
number, a 64-bit marker count, a 64-bit fingerprint, then for every marker the
six fields
\[
 (m_p,r_p,w_p,\eta_p,I_p,S_p).
\]
The first four entries are the invariant weight and phase-space variables
defined in equation (EWP2), \(I_p\) is the 64-bit persistent marker identifier,
and \(S_p\) is the 64-bit source-cell identifier that records the sample's
origin.  On reading, the driver rejects an incorrect tag or version, an
impossible count, invalid marker values, trailing bytes, or a recomputed
fingerprint that differs from the stored fingerprint.  If the requested path
already holds a different valid object, the writer refuses to replace it.

Ordinary continuation instead consumes the coordinate-specific atomic restart
in the output directory.  Before evolving, check finite positive invariant
weights, unique identifiers, support within the radial domain and
particle-shape stencil, convergence of the invariant-weight fixed point,
initial constraint residuals, ADM mass, compactness, binding diagnostics, and
quiet-pair identities where applicable.  A fixed-ADM or angular-momentum
family must retain the same authenticated base realization and be repeated
under marker, mesh, and timestep refinement.

\subsubsection{Task-specific input interface}

The two frontends are
\codesource{Einstein--Vlasov WPIC task frontends}
{\path{cpu/src/apps/einstein_vlasov/wpic_low_compactness.cpp};
 \path{cpu/src/apps/einstein_vlasov/wpic_black_hole.cpp}}
{Locks each executable to one low- or high-compactness problem contract before
 entering the shared relativistic marker driver}
\begin{lstlisting}[basicstyle=\ttfamily\scriptsize]
venus_cpu_einstein_vlasov_wpic_low_compactness
venus_cpu_einstein_vlasov_wpic_black_hole
\end{lstlisting}
The first accepts only the \code{low} scenario, while the second accepts only
\code{high}.
The shared options below retain the scenario prefixes to keep restart and
output contracts unambiguous.

\codesource{Einstein--Vlasov WPIC command-line interface}
{\path{cpu/src/drivers/einstein_vlasov_wpic.cpp}}
{Parses and validates the scenario-specific marker, grid, initial-condition,
 timestep, and output controls used by both frontends}
\begin{lstlisting}[basicstyle=\ttfamily\scriptsize]
# Scenario and execution
--scenario low|high                     (locked by executable)
--coordinates schwarzschild             (polar--areal module)
--expected-outcome black_hole|halo|any  (black_hole)
--output-dir DIR
--threads N                             (hardware concurrency)
--physical-ic PATH                      (immutable version-1 base IC)
--initialize-only                       (write/authenticate IC, then stop)
--resume                                (atomic checkpoints and CSV append)
--quick                                 (small diagnostic defaults)

# Low-compactness EV-versus-VP comparison
--low-markers N                         (16000)
--low-radial-cells N                    (192)
--low-t-end T                           (45)
--low-dt-max DT                         (0.06)
--low-cfl C                             (0.3)
--low-total-mass M                      (0.002)
--low-velocity-factor F                 (0.6)
--low-radial-importance-exponent P      (1)
--low-output-steps N                    (20)

# High-compactness Einstein--Vlasov problem
--high-initial-condition plummer|shell|energy_cutoff_shell (plummer)
--high-inner-boundary reflecting|origin_parity (origin_parity)
--high-radial-grid log|linear           (linear)
--high-markers N                        (20000)
--high-radial-cells N                   (320)
--high-t-end T                          (8)
--high-dt-max DT                        (0.008)
--high-gap-stop G                       (0.03)
--high-r-min R                          (0.01 input; origin parity resolves 0)
--high-r-max R                          (8)
--high-total-mass M                     (0.42)
--high-scale-radius A                   (1)
--high-velocity-factor F                (0.2)
--high-radial-infall W                  (0)
--high-infall-core-radius R             (0)
--high-radial-importance-exponent P     (1)
--high-output-steps N                   (10)
--high-target-initial-adm-mass M        (-1: keep base weights)
--high-specific-angular-momentum L      (-1: keep sampled values)
--high-angular-momentum-scale S         (1)
--high-sampling-realization UINT64      (0)

# Additional infalling-shell controls
--high-shell-center R                   (1.8)
--high-shell-width SIGMA_R              (0.18)
--high-shell-radial-importance-factor F (1)
--high-shell-radial-momentum W          (-0.62)
--high-shell-radial-momentum-dispersion SIGMA_W (0.025)
--high-shell-tangential-velocity-dispersion SIGMA_T (0.012)

# Bound energy-cutoff-shell controls
--high-energy-cutoff-amplitude C        (0.50)
--high-energy-cutoff-ell0 ELL0          (1.0)
--high-energy-cutoff-central-y Y0       (0.13)

# Shared weighted timestep tail
--timestep-omitted-mass-fraction F      (1e-5)
\end{lstlisting}
The \code{low} scenario evolves the same deterministic marker realization
with Einstein--Vlasov and Vlasov--Poisson, so differences measure relativistic
physics and discretization rather than different Monte Carlo noise.  The
shared driver also parses the horizon-penetrating coordinate-map and excision
controls, but those controls are explained only in the EF chapter, after that
geometry and its variables have been defined.

\noindent\textbf{Runnable high-compactness Schwarzschild example.}
The WPIC executables use command-line options rather than \path{params.ini}.
Build the task-specific high-compactness frontend and run it as follows:
\codesource{\code{venus_cpu_einstein_vlasov_wpic_black_hole}}
{\path{cpu/src/apps/einstein_vlasov/wpic_black_hole.cpp};
 \path{cpu/src/drivers/einstein_vlasov_wpic.cpp}}
{Selects the high-compactness problem contract; the shared driver performs its
 marker sampling, Schwarzschild metric solve, transport, stopping proxy,
 restart, and diagnostics}
\begin{lstlisting}[basicstyle=\ttfamily\scriptsize]
cd VENUS/cpu
make venus_cpu_einstein_vlasov_wpic_black_hole

./venus_cpu_einstein_vlasov_wpic_black_hole \
  --scenario high \
  --coordinates schwarzschild \
  --high-initial-condition plummer \
  --high-inner-boundary origin_parity \
  --high-radial-grid linear \
  --high-markers 128000 \
  --high-radial-cells 1024 \
  --high-r-max 4 \
  --high-t-end 20 \
  --high-dt-max 0.008 \
  --high-gap-stop 0.01 \
  --high-total-mass 0.6 \
  --high-scale-radius 1 \
  --high-velocity-factor 0.2 \
  --high-radial-infall 0.45 \
  --high-infall-core-radius 0.1 \
  --high-radial-importance-exponent 2 \
  --timestep-omitted-mass-fraction 1e-4 \
  --threads 32 \
  --output-dir output/ev_wpic_high
\end{lstlisting}
The executable name identifies the high-compactness task profile; it does not
certify that a black hole formed.  In particular, \code{--high-gap-stop} is a
Schwarzschild coordinate-patch stopping proxy.

\subsubsection{Output schemas and plotting}

The low task writes \path{low_compactness_history.csv} with this exact schema:
\codesource{Low-compactness EV--VP comparison output}
{\path{cpu/src/drivers/einstein_vlasov_wpic.cpp}}
{Writes relativistic and Newtonian mass, size, virial, density, and
 marker-difference diagnostics from the same initial realization}
\begin{lstlisting}[basicstyle=\ttfamily\scriptsize]
step,time,dt,ev_rest_mass,ev_adm_mass,ev_max_compactness,
ev_min_gap,ev_r10,ev_r50,ev_r90,ev_virial_ratio,
vp_mass,vp_r10,vp_r50,vp_r90,vp_virial_ratio,
density_l1_difference,particle_rms_radius_difference,
particle_rms_momentum_difference
\end{lstlisting}
Here \code{step} is the number of accepted time steps, \code{time} is the
current coordinate time, and \code{dt} is the most recently accepted time
increment; the initial row uses \code{dt=0}.  Prefix \code{ev_} denotes the
Einstein--Vlasov state and \code{vp_} the Vlasov--Poisson control.
\code{rest_mass} or \code{mass} is the sum of active marker weights,
\code{adm_mass} is the relativistic total gravitating mass,
\code{max_compactness} is \(\max_i(2m_i/r_i)\), and \code{min_gap} is
\(\min_i(1-2m_i/r_i)\).  The quantities \code{r10}, \code{r50}, and
\code{r90} are the radii containing 10, 50, and 90 percent of active marker
weight.  A \code{virial_ratio} is the Newtonian diagnostic \(2K/|W|\), where
\(K\) and \(W\) are the marker kinetic and spherical potential-energy sums;
it is a comparison aid, not a relativistic virial theorem.

The density difference is the normalized shell-mass distance
\[
 D_1=\frac{\sum_i|M_i^{\rm EV}-M_i^{\rm VP}|}
 {\tfrac12\sum_i(M_i^{\rm EV}+M_i^{\rm VP})}.
\]
The two particle differences are unweighted root-mean-square differences in
\(r_p\) and \(w_p\), evaluated marker by marker from the same initial realization.
The companion \path{low_compactness_profiles.csv} contains
\codesource{Low-compactness EV--VP radial-profile schema}
{\path{cpu/src/drivers/einstein_vlasov_wpic.cpp}}
{Writes the relativistic and Newtonian radial profile columns}
\begin{lstlisting}[basicstyle=\ttfamily\scriptsize]
step,time,r,ev_rho,ev_compactness,vp_rho,vp_compactness
\end{lstlisting}
where \code{r} is a shell-center radius, each \code{rho} is deposited mass or
energy density, and each \code{compactness} is \(2m(r)/r\) for the named
model.

The Schwarzschild high task writes \path{high_compactness_history.csv}:
\codesource{High-compactness Schwarzschild output}
{\path{cpu/src/drivers/einstein_vlasov_wpic.cpp}}
{Writes the coordinate-patch, mass, compactness, lapse, size, and virial
 diagnostics for every saved high-compactness state}
\begin{lstlisting}[basicstyle=\ttfamily\scriptsize]
step,time,dt,rest_mass,adm_mass,max_compactness,min_gap,
min_lapse,max_log_redshift,r10,r50,r90,virial_ratio,
anisotropic_jeans_relative_l1
\end{lstlisting}
The shared time, mass, compactness, radius, and virial columns have the
definitions above.  \code{min_lapse} is the minimum of \(e^\mu\) over the
saved metric faces (or centers if face values are unavailable), and
\code{max_log_redshift} is \(-\min\mu\).  The last column is the
volume-weighted relative \(L^1\) residual of the static anisotropic Jeans
balance
\[
 p_r'+(\rho+p_r)\mu_r-\frac{p_t-2p_r}{r}=0,
\]
where \(p_t\) is the sum of the two tangential pressures.  It is a late-time
stationarity check.  The calculation excludes the two derivative-boundary
cells, nonfinite cells, and cells below \(10^{-6}\) of the peak density.  A
value of \(-1\) means that fewer than three supported cells, or no nonzero
normalization, made this diagnostic unavailable.

The exact \path{high_compactness_profiles.csv} schema is
\codesource{High-compactness Schwarzschild radial-profile schema}
{\path{cpu/src/drivers/einstein_vlasov_wpic.cpp}}
{Writes center and face-aware metric, matter, and coordinate-map profile
 columns}
\begin{lstlisting}[basicstyle=\ttfamily\scriptsize]
step,time,r,rho,compactness,m,a,lambda,mu,lambda_t,
p_radial,p_tangential_sum,j,metric_sample_kind,metric_cell,
face_left_r,face_right_r,face_left_lambda,face_right_lambda,
center_mu_r,face_left_mu,face_right_mu,face_left_mu_r,
face_right_mu_r
\end{lstlisting}
The first fields identify the saved state and sample radius.  Here \code{rho}
is energy density, \code{m} is enclosed Misner--Sharp mass, and
\code{compactness}=\(2m/r\).  In this Schwarzschild file the historical column
name \code{a} means the gap \(1-2m/r=e^{-2\lambda}\); it is not introducing an
EF variable.  The next metric columns are \(\lambda\), \(\mu\), and
\(\lambda_\tau\).  The source member and CSV column \code{lambda_t}
retain their legacy names for this derivative.  The matter columns are
radial pressure, the sum of both
tangential pressures, and radial current.  \code{metric_sample_kind} is
\code{cell_center} for an ordinary row or \code{outer_face} for the exact outer
metric face saved for coordinate-map work; \code{metric_cell} is its zero-based
cell index.  The remaining fields give the radii and metric values at the two
faces bracketing a center, together with the center and face values of
\(\mu_r\).  Matter fields in the extra outer-face row are zero because that row
stores a metric boundary sample rather than a matter cell.

Restarts are separate from these readable tables.  The low task alternates
paired checkpoints
\path{low_einstein_vlasov_A.restart}/\path{low_einstein_vlasov_B.restart} and
\path{low_vlasov_poisson_A.restart}/\path{low_vlasov_poisson_B.restart}; the
atomically replaced \path{low_restart_manifest.txt} names the common active
slot, time, and step.  The high Schwarzschild task atomically replaces
\path{high_einstein_vlasov.restart}.  A resume requires the checkpoint time and
step to agree with the existing CSV tails before it appends any row.

The high task also writes
\path{fixed_family_initial_bound_diagnostics.csv}.  Its rest weight is the
sum of invariant marker weights, its binding energy is rest weight minus ADM
mass, and its binding fraction is that difference divided by rest weight.  The
four energy-at-infinity fields give the minimum, maximum, rest-weighted mean,
and rest-weight fraction at or above one for the instantaneous initial proxy
\(E_{\rm proxy}=e^\mu\epsilon\); this proxy is not conserved in a dynamical
metric.

Finally, \path{summary.json} records configuration, extrema, final diagnostics,
classification checks, and provenance fingerprints.  In this union schema,
\(-1\) in an optional time, radius, mass, coordinate-map, or energy-cutoff
field means ``not available or not observed,'' not a physical negative value.
The associated availability Boolean must be checked before using such a field.
For this Schwarzschild task, nonnegative \code{horizon_time} and
\code{first_trapped_time} record the gap-stop threshold time only;
\code{apparent_horizon_radius} remains \(-1\) because this solver does not
locate an apparent horizon.  The EF-only summary fields are defined in the EF
chapter and remain unavailable here.

Render both scenarios with:
\codesource{\path{render_weighted_einstein_vlasov_pic_profile_gifs.py}}
{\path{cpu/tools/render_weighted_einstein_vlasov_pic_profile_gifs.py}}
{Reads the documented low- and high-profile schemas and renders their profile
 animations and preview products}
\begin{lstlisting}[basicstyle=\ttfamily\scriptsize]
python tools/render_weighted_einstein_vlasov_pic_profile_gifs.py \
  --low-profiles LOW_RUN/low_compactness_profiles.csv \
  --high-profiles HIGH_RUN/high_compactness_profiles.csv \
  --output-dir plots/ev_wpic
\end{lstlisting}
Here \code{LOW_RUN} and \code{HIGH_RUN} stand for the two output directories
chosen with \code{--output-dir}.  The plotting script validates the required
named columns, uses saved times without temporal interpolation, and writes
separate low- and high-profile GIF animations, up to three PNG previews for
each, and exact frame-time sidecars.

For a low-compactness comparison, use the low frontend and monitor the EV--VP
density and marker differences alongside ADM-mass conservation.  For collapse,
refine marker count, radial mesh, timestep controls, and the omitted-mass
fraction together.  The stopping gap and maximum compactness diagnose the
approach to \(2m/r=1\); they must not be interpreted as evolution through an
apparent horizon.

For a paired collapse, the maintained three-panel animation is
\path{cpu/tools/render_weighted_einstein_vlasov_coordinate_pair_gif.py}.  This
paragraph previews only the notation needed to explain that plot; the EF
chapter derives it.  Ingoing EF coordinates use
\[
 \dd s^2=-a(v,r)b(v,r)^2\dd v^2+2b(v,r)\dd v\dd r
          +r^2\dd\Omega^2,
\]
where \(v\) is advanced (ingoing-null) time, \(a=1-2m/r\) is the radial null
factor, and the positive function \(b\) fixes the normalization of \(v\).
Here and below SC abbreviates the Schwarzschild description.
Before the SC chart ends, the renderer does not pair curves merely because
their two coordinate clocks show the same number.  At every displayed radius
it solves the eikonal map---the equation that follows the same ingoing light
surface through the two coordinate systems---
\[
 \partial_r v=e^{\lambda-\mu}\partial_\tau v,
 \qquad v(\tau,R_B)=\tau+v_{\rm offset},
\]
where \(\tau\) is Schwarzschild time, \(R_B\) is the authenticated outer vacuum
boundary, and \(v_{\rm offset}\) aligns the two clocks there.  It checks the
raw \(n\)-versus-\(2n\) map-integration refinements, the inverse map, and their
composition, and compares SC at \((\tau,r)\) with EF at
\((v(\tau,r),r)\).  The panels show
\(\log_{10}\rho\) versus \(\log_{10}r\),
\(\dd\log\rho/\dd\log r\) versus \(\log_{10}r\), and \(2m/r\) versus the
linear radius \(r\).  Outside a trapped region the EF density shown in the
overlay is
\[
 \rho_{\rm EF}^{(\mathrm{SC\ observer})}=\frac{T_{vv}}{a b^2}.
\]
Here \(T_{vv}\) is the covariant stress--energy component with both indices in
the advanced-time direction; division by \(ab^2\) converts it to the energy
density measured by the corresponding static SC observer.
Where \(a\leq0\), that static observer no longer exists.  Density and its
slope are therefore masked there instead of being silently replaced by a
different quantity; compactness remains visible across the full saved EF
slice.  After the Schwarzschild chart ends, clearly labelled EF-only frames
show horizon and exterior evolution.  For example:
\codesource{\path{render_weighted_einstein_vlasov_coordinate_pair_gif.py}}
{\path{cpu/tools/render_weighted_einstein_vlasov_coordinate_pair_gif.py}}
{Authenticates the paired inputs, matches the same spacetime events through
 the solved eikonal map, and publishes the comparison animation and sidecars}
\begin{lstlisting}[basicstyle=\ttfamily\scriptsize]
python tools/render_weighted_einstein_vlasov_coordinate_pair_gif.py \
  --schwarzschild RUN/sc --ef RUN/ef \
  --manifest RUN/common_ingress_null.bin.manifest.json \
  --eikonal-substeps "$VENUS_EIKONAL_SUBSTEPS" \
  --output RUN/sc_ef_black_hole_profiles.gif
\end{lstlisting}
In this template, \code{RUN} is the parent directory containing the paired
\path{sc} and \path{ef} run directories.  The shell variable
\code{VENUS_EIKONAL_SUBSTEPS} is the positive map-integration refinement count
chosen for the authenticated comparison.
The GIF is published atomically with three PNG previews, an exact frame-clock
CSV, and a JSON record of input hashes, map checks, panel definitions, and
fixed axis limits.
